% =========================================================================
% PRÁCTICA 9: CONTROLADOR DE MATRIZ DE LEDS 8x8
% =========================================================================
\section{Práctica 9: Controlador de Matriz de LEDs 8x8}

\subsection{Descripción General}
Este desarrollo se centró en el diseño y programación de una interfaz visual dinámica mediante una matriz de diodos emisores de luz organizada en un arreglo ortogonal de $8\times8$[cite: 44]. Con la finalidad de superar las restricciones de pines físicas inherentes al encapsulado del microcontrolador, se implementó un algoritmo basado en la técnica de multiplexación por barrido secuencial de filas[cite: 44, 51]. La persistencia de visión del ojo humano unifica las ráfagas discretas de encendido configuradas a $2\text{ ms}$, proyectando figuras complejas e interactivas sin parpadeos perceptibles[cite: 47, 51].

\subsection{Componentes Utilizados}
\begin{itemize}
    \item \textbf{Unidad Central:} 1 Microcontrolador PIC16F877A, 1 Módulo programador PICkit 3[cite: 48, 49].
    \item \textbf{Unidad de Visualización:} 1 Matriz de LEDs común de $8\times8$[cite: 48].
    \item \textbf{Generador de Reloj:} 1 Cristal de cuarzo de $4\text{ MHz}$, 2 Capacitores de $22\text{ pF}$[cite: 48].
    \item \textbf{Interfaz de Entrada:} 1 Interruptor codificador múltiple (\textit{DIP Switch})[cite: 47].
    \item \textbf{Red Estabilizadora:} 4 Resistencias de $220\ \Omega$, 1 Resistencia de $1\text{ k}\Omega$, 1 Pulsador[cite: 48].
    \item \textbf{Interconexión:} Placa de pruebas (\textit{Protoboard}) y Jumpers de conexión rápida[cite: 48].
\end{itemize}

\subsection{Entorno de Desarrollo Software}
\begin{itemize}
    \item \textbf{IDE:} MPLAB X IDE v6.25[cite: 50].
    \item \textbf{Compilador:} MPLAB XC8 (Compiler Toolchain)[cite: 46].
    \item \textbf{Simulación Electrónica:} Proteus VSM 8[cite: 50].
\end{itemize}

\subsection{Principios de Operación y Control}
La lógica del sistema explota el almacenamiento de mapas de bits indexados en arreglos globales unidimensionales de tipo \texttt{char}[cite: 48, 52]. El Puerto D actúa de forma síncrona como el bus selector de filas (operando en cátodo común, activas en nivel bajo), mientras que el Puerto B es el encargado de inyectar los datos lógicos en las columnas (ánodos, activos en nivel alto)[cite: 47, 51].

El ciclo de refresco (\textit{refresh rate}) es ejecutado de forma continua por la rutina \texttt{renderizar\_matriz()}[cite: 53]. Esta rutina recorre mediante un bucle indexado las ocho filas, desplazando un cero lógico (\texttt{\~{}(1 << i)}) para habilitar selectivamente una línea a la vez mientras escribe la palabra binaria homóloga guardada en el búfer[cite: 65, 66]. Para evitar falsos disparos por ruido electromecánico en los interruptores del Puerto A, se diseñó una función de lectura con doble verificación temporal (\textit{debounce} por software)[cite: 70].

\subsection{Código de Control en C}
\begin{lstlisting}[language=C, caption={Firmware para el control multiplexado de la matriz 8x8}, label={code:practica9}]
#pragma config FOSC = XT       // Reloj oscilador por cristal de cuarzo
#pragma config WDTE = OFF      // Deshabilitacion del temporizador perro guardian
#pragma config PWRTE = ON      // Retraso temporizado al encender
#pragma config BOREN = ON      // Monitoreo de caida de tension activo
#pragma config LVP = OFF       // Programacion de bajo voltaje inhabilitada
#pragma config CPD = OFF       // Sin proteccion de lectura en EEPROM
#pragma config CP = OFF        // Sin proteccion de lectura en codigo

#define _XTAL_FREQ 4000000     // Frecuencia base de 4MHz

#include <xc.h>

// Declaracion de prototipos funcionales
void inicializar_sistema(void);
void renderizar_matriz(void);
unsigned char capturar_entrada(void);

// Iconografia cargada en matrices de datos estaticos
void cargar_vacio(void);
void cargar_creeper(void);
void cargar_corazon(void);
void cargar_simbolo(void);
void cargar_personalizado(void);
void cargar_test(void);

// Buffer dinamico de visualizacion
unsigned char buffer_grafico[8] = {0x00, 0x00, 0x00, 0x00, 0x00, 0x00, 0x00, 0x00};

void main(void) {
    unsigned char token_seleccion;
    inicializar_sistema();
    
    while(1) {
        token_seleccion = capturar_entrada(); // Captura de estados en DIP switch
        
        switch(token_seleccion) {
            case 0: cargar_vacio();         break; // Binario 000
            case 1: cargar_creeper();       break; // Binario 001
            case 2: cargar_corazon();       break; // Binario 010
            case 3: cargar_simbolo();       break; // Binario 011
            case 4: cargar_personalizado(); break; // Binario 100
            case 7: cargar_test();          break; // Binario 111
            default: cargar_personalizado(); break;
        }
        
        renderizar_matriz(); // Proyeccion multiplexada continua
    }
}

void inicializar_sistema(void) {
    ADCON1 = 0x06;  // Desconexion del bloque analogico, Puerto A digital
    TRISA  = 0xFF;  // Puerto A integro mapeado como lineas de entrada
    TRISB  = 0x00;  // Puerto B mapeado como bus de salida (Anodos de columna)
    TRISD  = 0x00;  // Puerto D mapeado como bus de salida (Catodos de fila)
    PORTB  = 0x00;  // Limpieza inicial de columnas
    PORTD  = 0xFF;  // Lineas de fila desactivadas (1 logico apaga el catodo)
}

void renderizar_matriz(void) {
    for(char i = 0; i < 8; i++) {
        PORTD = (unsigned char)~(1 << i); // Inyeccion del cero secuencial en filas
        PORTB = buffer_grafico[i];         // Transmision del patron de columnas
        __delay_ms(2);                     // Ventana temporal critica de persistencia
    }
}

unsigned char capturar_entrada(void) {
    unsigned char lectura_Alfa, lectura_Beta;
    do {
        lectura_Alfa = PORTA & 0x07; // Filtrado de las 3 lineas bajas del puerto
        __delay_ms(10);
        lectura_Beta = PORTA & 0x07;
    } while(lectura_Alfa != lectura_Beta); // Validacion de estabilidad estructural
    
    return lectura_Alfa;
}

void cargar_creeper(void) {
    unsigned char mapa[8] = {0x00, 0x00, 0x66, 0x66, 0x18, 0x3C, 0x3C, 0x24};
    for(char i = 0; i < 8; i++) buffer_grafico[i] = mapa[i];
}

void cargar_vacio(void) {
    for(char i = 0; i < 8; i++) buffer_grafico[i] = 0x00;
}

void cargar_corazon(void) {
    unsigned char mapa[8] = {0x00, 0x66, 0xFF, 0xFF, 0x7E, 0x3C, 0x18, 0x00};
    for(char i = 0; i < 8; i++) buffer_grafico[i] = mapa[i];
}

void cargar_simbolo(void) {
    unsigned char mapa[8] = {0x18, 0x3C, 0x7E, 0xDB, 0x99, 0x18, 0x18, 0x18};
    for(char i = 0; i < 8; i++) buffer_grafico[i] = mapa[i];
}

void cargar_personalizado(void) {
    unsigned char mapa[8] = {0xFF, 0x81, 0xBD, 0xA5, 0xA5, 0xBD, 0x81, 0xFF};
    for(char i = 0; i < 8; i++) buffer_grafico[i] = mapa[i];
}

void cargar_test(void) {
    for(char i = 0; i < 8; i++) buffer_grafico[i] = 0xFF;
}
\end{lstlisting}

\subsection{Conclusiones}
La resolución de este proyecto demostró la efectividad de la multiplexación dinámica por software como solución óptima para extender los recursos físicos de entrada y salida del PIC16F877A[cite: 75, 79]. Reducir la demanda de pines lógicos de 64 líneas independientes a únicamente dos puertos estándar optimiza drásticamente la topología del circuito impreso o de prueba[cite: 79]. Adicionalmente, se concluye que el control temporal preciso en el refresco secuencial ($2\text{ ms}$) es el pilar del diseño, dado que una ventana mayor induciría parpadeos indeseados, rompiendo la persistencia de visión necesaria en la interfaz[cite: 78, 79].
